\documentclass[10pt]{book}
\usepackage{ctex}
\usepackage{amsmath, amsthm, amssymb, amsfonts}
\usepackage{thmtools}
\usepackage[dvipsnames]{xcolor}
\usepackage{tcolorbox}
\tcbuselibrary{skins,breakable,theorems}
\usepackage{hyperref}
\usepackage{fontspec}
\usepackage{tikz}
\usepackage{fancyhdr}
\pagestyle{fancy}
\usetikzlibrary{arrows.meta, positioning}

% -------------------------------
% 全局正文字体（英文 + 中文） + 伪粗体
% -------------------------------
\setmainfont{Times New Roman}
\setCJKmainfont{SimSun}[AutoFakeBold={3}]
\setCJKfamilyfont{theoremCJK}{FangSong}[AutoFakeBold={2.5}]
\setCJKfamilyfont{proofCJK}{KaiTi}[AutoFakeBold={2.5}]
\newcommand{\theoremFont}{\CJKfamily{theoremCJK}}

% 证明字体
\newcommand{\proofFont}{\CJKfamily{proofCJK}}

% -------------------------------
% 鲜明蓝色系配色方案
% -------------------------------
\definecolor{LogicBlue}{RGB}{0, 102, 204}
\definecolor{LogicLightBlue}{RGB}{102, 204, 255}
\definecolor{LogicDarkBlue}{RGB}{0, 51, 153}
\definecolor{LogicAccentBlue}{RGB}{51, 153, 255}

% -------------------------------
% 逻辑符号快捷命令
% -------------------------------
\newcommand{\landand}{\land}
\newcommand{\loror}{\lor}
\newcommand{\limplies}{\to}
\newcommand{\liff}{\leftrightarrow}
\newcommand{\entails}{\models}
\newcommand{\turnstile}{\vdash}

% -------------------------------
% 练习（无编号）
% -------------------------------
\NewDocumentEnvironment{exercise}{}{\par\noindent{\bfseries\theoremFont 练习\quad}}{\par}

% -------------------------------
% 彩色 tcolorbox 环境 —— 所有共用同一计数器（以 mydefinition 的计数器为主）
% -------------------------------
% 第一个环境：正常带 section 内编号
\newtcbtheorem[number within=section]{mydefinition}{定义}{
	breakable, enhanced,
	colback=white, colframe=LogicBlue,
	arc=4pt, boxrule=0.8pt, leftrule=4pt,
	fonttitle=\bfseries\theoremFont,
	fontupper=\normalsize\CJKfamily{theoremCJK},
	left=16pt, right=12pt, top=10pt, bottom=10pt,
	before skip=12pt, after skip=12pt,
	boxsep=4pt, left skip=16pt,
}{df}

% 后续环境：借用 mydefinition 的计数器（tcb@cnt@df）
\newtcbtheorem[use counter from=mydefinition]{mytheorem}{定理}{
	breakable, enhanced,
	colback=white, colframe=LogicBlue,
	arc=4pt, boxrule=0.8pt, leftrule=4pt,
	fonttitle=\bfseries\theoremFont,
	fontupper=\normalsize\CJKfamily{theoremCJK},
	left=16pt, right=12pt, top=10pt, bottom=10pt,
	before skip=12pt, after skip=12pt,
	boxsep=4pt, left skip=16pt,
}{th}

\newtcbtheorem[use counter from=mydefinition]{mycorollary}{推论}{
	breakable, enhanced,
	colback=white, colframe=LogicLightBlue,
	arc=4pt, boxrule=0.8pt, leftrule=4pt,
	fonttitle=\bfseries\theoremFont,
	fontupper=\normalsize\CJKfamily{theoremCJK},
	left=16pt, right=12pt, top=10pt, bottom=10pt,
	before skip=12pt, after skip=12pt,
	boxsep=4pt, left skip=16pt,
}{cor}

\newtcbtheorem[use counter from=mydefinition]{myaxiom}{公理}{
	breakable, enhanced,
	colback=white, colframe=LogicDarkBlue,
	arc=4pt, boxrule=0.8pt, leftrule=4pt,
	fonttitle=\bfseries\theoremFont,
	fontupper=\normalsize\CJKfamily{theoremCJK},
	left=16pt, right=12pt, top=10pt, bottom=10pt,
	before skip=12pt, after skip=12pt,
	boxsep=4pt, left skip=16pt,
}{ax}

\newtcbtheorem[use counter from=mydefinition]{myinferencerule}{推理规则}{
	breakable, enhanced,
	colback=white, colframe=LogicAccentBlue,
	arc=4pt, boxrule=0.8pt, leftrule=4pt,
	fonttitle=\bfseries\theoremFont,
	fontupper=\normalsize\CJKfamily{theoremCJK},
	left=16pt, right=12pt, top=10pt, bottom=10pt,
	before skip=12pt, after skip=12pt,
	boxsep=4pt, left skip=16pt,
}{ir}

\newtcbtheorem[use counter from=mydefinition]{myexample}{范例}{
	breakable, enhanced,
	colback=white, colframe=LogicLightBlue,
	arc=4pt, boxrule=0.8pt, leftrule=4pt,
	fonttitle=\bfseries\theoremFont,
	fontupper=\normalsize\CJKfamily{theoremCJK},
	left=16pt, right=12pt, top=10pt, bottom=10pt,
	before skip=12pt, after skip=12pt,
	boxsep=4pt, left skip=16pt,
}{ex}

\newtcbtheorem[use counter from=mydefinition]{myremark}{注记}{
	breakable, enhanced,
	colback=white, colframe=LogicDarkBlue,
	arc=4pt, boxrule=0.8pt, leftrule=4pt,
	fonttitle=\bfseries\theoremFont,
	fontupper=\normalsize\CJKfamily{theoremCJK},
	left=16pt, right=12pt, top=10pt, bottom=10pt,
	before skip=12pt, after skip=12pt,
	boxsep=4pt, left skip=16pt,
}{rm}

% -------------------------------
% 证明环境（楷体，无边框，不参与编号）
% -------------------------------
\RenewDocumentEnvironment{proof}{}{\par\noindent{\bfseries\proofFont 证明.\quad}}{\hfill$\square$\par\bigskip}

% -------------------------------
% 真值表环境（实用版）
% -------------------------------
\newenvironment{truthtable}[1][c|c|c|c|c|c]{%
	\par\noindent\centering
	\begin{tabular}{|#1|}\hline
		\addlinespace[0.4ex]%
	}{%
		\hline\end{tabular}\par\bigskip
}

% -------------------------------
% 页面样式
% -------------------------------
\fancyhf{}
\fancyfoot[C]{\thepage}
\renewcommand{\headrulewidth}{0pt}
\renewcommand{\footrulewidth}{0pt}

\usepackage[bottom=1.5cm]{geometry}
\setlength{\footskip}{1cm}